\item 
\textbf{HydroSense Tech }    \hfill     
\hfill  \textit{Sep 2017 -- Present (Part Time, remote); 
\item  Co-founder / Member of Technical Staff\hfill May 2025 -- Sep 2025 (Full Time, Singapore)}
\vspace*{-3pt}
\begin{itemize}[leftmargin=0.2in, topsep=1pt]
  \item Co-founded an IoT startup and designed evaluation pipelines enabling models to self-improve through end-user interactions, building features and data flows that capture human feedback to drive continuous optimization of user value and engagement metrics.
  \item Built RLHF tooling to gather and incorporate end-user preferences and feedback, enabling models to iteratively self-improve based on real-world usage patterns with scalable data pipelines and preference-based learning systems.
  \item Developed data pipelines and AI evaluation frameworks to track model performance, safety, user satisfaction, and metrics like retention and engagement from an end-user lens, implementing custom benchmarks and A/B testing infrastructure.
  \item Worked with product and model training teams to optimize and productionize AI-powered features, scaling from prototype to global deployment with robust evaluation methodologies and data-driven insights for alignment and user-centric improvement.
  \item \textbf{Patents:} Co-inventor on 3 granted patents covering rainfall sensing hardware, ML-based drift compensation, and intelligent embedded calibration algorithms for industrial monitoring applications.
\end{itemize}

\vspace*{-1pt}
\item \textbf{Seno Medical}\\
\textit{Member of Technical Staff} \hfill Buffalo, NY \quad \textit{May 2024 -- Aug 2024}
\vspace*{-3pt}
\begin{itemize}[leftmargin=0.2in, topsep=1pt]
  \item Collaborated with domain experts to build a \textbf{2k+ QA} dataset from technical documents; fine-tuned a \textbf{LLaMA-based model (LoRA)} and implemented retrieval-augmented generation (RAG) system with evaluation tools and frameworks for production deployment.
  \item Designed an end-to-end evaluation harness with label schema (factuality, coverage, style), automatic metrics using RAGAS and TruLens, and Python error-analysis scripts; reduced hallucination rate from \textbf{30\% to 15\%} and increased \textbf{Recall@5 to 90\%} through iterative feedback loops.
  \item Built production-grade Python services for AI and ML algorithms, implementing large language model inference via \textbf{FastAPI} and containerized with \textbf{Docker}, ensuring scalable and high-performance AI infrastructure with comprehensive monitoring and observability.
  \item Set up Git-based workflows and production-grade observability with structured logs and monitoring dashboards, shipping reproducible Docker images with deployment runbooks and maintaining high code quality standards for rapid iteration cycles.
\end{itemize}

\vspace*{-1pt}
\item \textbf{Roswell Park Comprehensive Cancer Center}\\
\textit{Research Scientist} \hfill Buffalo, NY \quad \textit{May 2023 -- Aug 2023}
\vspace*{-3pt}
\begin{itemize}[leftmargin=0.2in, topsep=1pt]
  \item Designed and implemented a multi-modal data pipeline aligning imaging with structured text reports, including de-identification, schema normalization, and automated QC, to produce clean image–text pairs for generative modeling and downstream training.
  \item Built dataset generation jobs with filtering, sampling, augmentation, and metadata tagging, creating balanced image–report datasets with vector database integration for conditional generation and enabling fast iteration over model variants.
  \item Implemented a domain-specific latent diffusion model (Stable Diffusion–style, DDIM/DPMS sampling) trained on \textbf{AWS}, integrating LLM orchestration for conditional image–to–text report generation and synthetic data augmentation with evaluation tracking.
  \item Automated GPU training and experiment tracking with standardized configs, observability pipelines, and comparison scripts, enabling fast iteration over model variants and dataset settings on distributed AWS infrastructure with data-driven insights.
\end{itemize}
